
\documentclass{article}
\usepackage[utf8]{inputenc}
\usepackage{amsmath}
\usepackage{amsfonts}
\usepackage{amssymb}
\usepackage{geometry}
\geometry{a4paper, margin=15mm}

\title{2025 MIT Science Bowl High School Invitational - Round 16}
\author{}
\date{}

\begin{document}

\maketitle

\section*{Round 16}

\subsection*{TOSS-UP 1}
\textbf{CHEMISTRY - Multiple Choice} Which of the following is the mathematical origin of planar nodes in a hydrogen-like atom's atomic orbitals? \\
W) Roots of radial wavefunction \\
X) Critical points of radial wavefunction \\
Y) Roots of angular wavefunction \\
Z) Critical points of angular wavefunction \\
\textbf{ANSWER: Y) ROOTS OF ANGULAR WAVEFUNCTION}

\subsection*{VISUAL BONUS 1}
\textbf{CHEMISTRY - Short Answer} Your next bonus is visual. Shown in the image is the unit cell of a hexagonal binary metal oxide. Answer the following two questions about this metal oxide:
1) What is the oxidation state of the metal in this oxide?
2) Identify all of the following three symmetries that exist in this mineral's crystal structure: a) Twofold rotation; b) Threefold rotation; c) Inversion through a point. \\
\textbf{ANSWER: 1) +2; 2) B ONLY}

\subsection*{TOSS-UP 2}
\textbf{EARTH AND SPACE - Multiple Choice} Which of the following features is NOT depositional in origin? \\
W) Alluvial Fan \\
X) Drumlin \\
Y) Roche Moutonnee \\
Z) Tombolo \\
\textbf{ANSWER: Y) ROCHE MOUTONNEE}

\subsection*{VISUAL BONUS 2}
\textbf{EARTH AND SPACE - Short Answer} Your next bonus is visual. Answer the following two questions regarding the change in rainfall patterns on the global map.
1) Which phase of the ENSO cycle is depicted on this map?
2) What effect does this weather pattern have on hurricane activity in the Atlantic ocean? \\
\textbf{ANSWER: 1) EL NINO; 2) DECREASE}

\subsection*{TOSS-UP 3}
\textbf{PHYSICS - Multiple Choice} Which of the following is an example of a unitary operator in quantum mechanics? \\
W) Translation \\
X) Position \\
Y) Spin \\
Z) Angular momentum \\
\textbf{ANSWER: W) TRANSLATION}

\subsection*{VISUAL BONUS 3}
\textbf{PHYSICS - Short Answer} Your next bonus is visual. Shown in this image is a system of simple machines. Assume all components are frictionless and massless. Knowing that the floor and ceiling are immovable, what is the system's ideal mechanical advantage? \\
\textbf{ANSWER: 4/3}

\subsection*{TOSS-UP 4}
\textbf{ENERGY - Short Answer} Researchers at the Broad institute have been studying Cuproptosis, a copper induced form of apoptosis. Cuproptosis occurs when copper destroys the iron sulfur clusters found in which group of mitochondrial membrane proteins? \\
\textbf{ANSWER: CYTOCHROMES}

\subsection*{BONUS 4}
\textbf{ENERGY - Short Answer} The Furst Lab used polyphenols to protect microbes responsible for producing fertilizers. Identify all of the following three biomolecules that are polyphenols: 1) Lignin; 2) Tannic acid; 3) Terpene. \\
\textbf{ANSWER: 1 AND 2}

\subsection*{TOSS-UP 5}
\textbf{MATH - Short Answer} Let $f(x)=a^{x}+b^{x}$, where a and b are positive integers. Given that $f(1)=9$ and $f(3)=351$, find the value of $f(2)$. \\
\textbf{ANSWER: 53}

\subsection*{VISUAL BONUS 5}
\textbf{MATH - Short Answer} Your next bonus is visual. Katherine and Selena are playing the game "Kokushi Challenge," which uses a set of 52 tiles with four identical copies of each tile. However, Katherine has identified that there is a 50\% chance that Selena has tampered with the set, painting one of the four "haku" tiles into another tile. He shuffles the set randomly and draws three tiles at random, and finds that they are all "haku" tiles. With this knowledge, what is the new probability that Selena has cheated? \\
\textbf{ANSWER: 1/5}

\subsection*{TOSS-UP 6}
\textbf{BIOLOGY - Multiple Choice} The presence of antinuclear antibodies in serum most likely indicates which of the following diseases? \\
W) Systemic lupus erythematosus \\
X) Acute lymphoblastic leukemia \\
Y) Cushing syndrome \\
Z) Hemolytic anemia \\
\textbf{ANSWER: W) SYSTEMIC LUPUS ERYTHEMATOSUS}

\subsection*{VISUAL BONUS 6}
\textbf{BIOLOGY - Short Answer} Your next bonus is visual. Shown in the image is a cross section of a specific type of leaf. Answer the following two questions:
1) Is this plant a xerophyte, hydrophyte, or mesophyte?
2) Identify the structure labeled. \\
\textbf{ANSWER: 1) HYDROPHYTE; 2) SCLERID}

\subsection*{TOSS-UP 7}
\textbf{PHYSICS - Multiple Choice} Which of the following changes would increase the photocurrent produced by shining light on a semiconductor? \\
W) Increasing the wavelength of light \\
X) Increasing the intensity of light \\
Y) Decreasing the wavelength of light \\
Z) Decreasing the intensity of light \\
\textbf{ANSWER: X) INCREASING THE INTENSITY OF LIGHT}

\subsection*{BONUS 7}
\textbf{PHYSICS - Short Answer} One metal wire runs along the x-axis and another wire runs along the y-axis, each carrying a current of 1 ampere in the positive direction. Find the ratio between the magnetic field strength at the point (2,1) and the magnetic field strength at (3,-1). \\
\textbf{ANSWER: 3/8 (ACCEPT: 0.375)}

\subsection*{TOSS-UP 8}
\textbf{CHEMISTRY - Multiple Choice} Cisplatin is a neutral platinum complex containing two chloride ligands and two ammonia ligands. Which of the following is not true about cisplatin? \\
W) Cisplatin is optically active \\
X) Cisplatin is diamagnetic \\
Y) Cisplatin is planar \\
Z) Cisplatin is polar \\
\textbf{ANSWER: W) CISPLATIN IS OPTICALLY ACTIVE}

\subsection*{BONUS 8}
\textbf{CHEMISTRY - Short Answer} Rank the following three substituents by increasing activation of an aryl ring towards electrophilic aromatic substitution: 1) Tribromomethyl; 2) Tert-butyl; 3) Methyl. \\
\textbf{ANSWER: 1, 3, 2}

\subsection*{TOSS-UP 9}
\textbf{MATH - Short Answer} Suppose a twice-differentiable function f satisfies $f(0)=0$, $f(1)=10$, and $f^{\prime\prime}(x)\ge0$ for all $x\in[0,1]$. Identify all of the following three options that are possible values for $f(0.5)$: 1) 2.5; 2) 5; 3) 7.5 \\
\textbf{ANSWER: 1,2 (ACCEPT: ALL BUT 3)}

\subsection*{BONUS 9}
\textbf{MATH - Short Answer} Brian calculates the expected value of the taxicab distance in n-dimensional space between a uniformly randomly chosen point and its nearest lattice point. This quantity is directly proportional to n raised to what power? \\
\textbf{ANSWER: 1 (ACCEPT: FIRST)}

\subsection*{TOSS-UP 10}
\textbf{ENERGY - Short Answer} The Kulik Group developed strong cross-linked molecular nets by substituting carbon atoms with atoms of what chemically similar element? \\
\textbf{ANSWER: SILICON}

\subsection*{BONUS 10}
\textbf{ENERGY - Multiple Choice} The Kulik Group is studying enzymes that catalyze biomolecular SN2 reactions. When bound to the substrate, the active sites of these enzymes most closely approximate which of the following molecular geometries? \\
W) Trigonal Planar \\
X) Tetrahedral \\
Y) Trigonal Pyramidal \\
Z) Trigonal Bipyramidal \\
\textbf{ANSWER: Z) TRIGONAL BIPYRAMIDAL}

\subsection*{TOSS-UP 11}
\textbf{BIOLOGY - Multiple Choice} The dense bodies of smooth muscle are analagous to which of the following structures in skeletal muscle? \\
W) Z disc \\
X) Thin filament \\
Y) Thick filament \\
Z) A band \\
\textbf{ANSWER: W) Z DISC}

\subsection*{VISUAL BONUS 11}
\textbf{BIOLOGY - Short Answer} Your next bonus is visual. Shown in the image is a Lotka-Volterra competition model used for representing competition between two species A and B. The lines represent isoclines at which a species has no net change in population, with blue representing species A and red representing species B. Answer the following two questions:
1) To the nearest ten, what are the coordinates corresponding to all equilibria in the model?
2) If the initial populations of A and B are 75 and 50, respectively, which species is more likely to dominate in competition? \\
\textbf{ANSWER: 1) (0,0), (0,100), (90,10), (150,0) IN ANY ORDER; 2) SPECIES B}

\subsection*{TOSS-UP 12}
\textbf{EARTH AND SPACE - Short Answer} Order the following three types of plagioclase feldspar from most sodium-rich to most calcium-rich: 1) Albite; 2) Anorthite; 3) Labradorite. \\
\textbf{ANSWER: 1, 3, 2}

\subsection*{BONUS 12}
\textbf{EARTH AND SPACE - Multiple Choice} Mr. Maine observes concentric rings of color change, or pleochroic halos, surrounding tiny inclusions of zircon in a biotite crystal. What do these halos indicate regarding the history of the zircon inclusions? \\
W) They were introduced during deformation, causing dislocation halos \\
X) They carried water and hydrated surrounding minerals \\
Y) They acted as nucleation sites that caused some biotite to crystallize first \\
Z) They contained radioactive elements that emitted alpha particles over time \\
\textbf{ANSWER: Z) THEY CONTAINED RADIOACTIVE ELEMENTS THAT EMITTED ALPHA PARTICLES OVER TIME}

\subsection*{TOSS-UP 13}
\textbf{PHYSICS - Short Answer} In general relativity, what tensor represents the source of the gravitational field? \\
\textbf{ANSWER: STRESS-ENERGY TENSOR}

\subsection*{BONUS 13}
\textbf{PHYSICS - Short Answer} The north pole of a small bar magnet is moving at constant velocity toward a very small loop of copper wire, such that the axis of the magnet is perpendicular to the plane of the loop. When the magnet is 2 meters from the wire, 1 microampere of current is induced in the wire. To the nearest microampere, how much current is induced at a separation distance of 1 meter? \\
\textbf{ANSWER: 16}

\subsection*{TOSS-UP 14}
\textbf{CHEMISTRY - Short Answer} Rank the following three carboxylic acid derivatives by increasing stretching frequency of their carbonyl double bond: 1) Amide; 2) Ester; 3) Acid bromide. \\
\textbf{ANSWER: 1, 2, 3}

\subsection*{VISUAL BONUS 14}
\textbf{CHEMISTRY - Short Answer} Your next bonus is visual. Shown in the image is a homogenous catalyst used in olefin metathesis. Answer the following two questions about this image:
1) What is the oxidation state of the catalyst's molybdenum metal center?
2) What is the electron count of molybdenum in this complex? \\
\textbf{ANSWER: 1) +6; 2) 12}

\subsection*{TOSS-UP 15}
\textbf{BIOLOGY - Short Answer} Identify all of the following three types of cellular junctions that use connexins: 1) Gap junction; 2) Adherens junction; 3) Desmosomes. \\
\textbf{ANSWER: 1 ONLY}

\subsection*{BONUS 15}
\textbf{BIOLOGY - Multiple Choice} During an enzyme assay, initial product formation is rapid before leveling off into a slower, steady-state rate. This burst-phase kinetic behavior is typically observed when the enzyme forms a covalent intermediate during catalysis. Which of the following enzymes most likely exhibits this behavior? \\
W) Amylase \\
X) Catalase \\
Y) Chymotrypsin \\
Z) Carbonic anhydrase \\
\textbf{ANSWER: Y) CHYMOTRYPSIN}

\subsection*{TOSS-UP 16}
\textbf{EARTH AND SPACE - Multiple Choice} Which of the following best explains the presence of acoustic peaks in the CMB power spectrum? \\
W) Interference between redshifted gravitational waves \\
X) Standing pressure waves in the photon-baryon plasma \\
Y) Variations in neutrino mass density at recombination \\
Z) Phase transition from quark-gluon plasma \\
\textbf{ANSWER: X) STANDING PRESSURE WAVES IN THE PHOTON-BARYON PLASMA}

\subsection*{VISUAL BONUS 16}
\textbf{EARTH AND SPACE - Short Answer} Your next bonus is visual. Answer the following two questions about this histogram of data collected from Cassini's flybys of particle-rich jets:
1) From which celestial object did those jets originate?
2) Which molecules are represented by the peaks labeled A and B, respectively? \\
\textbf{ANSWER: 1) ENCELADUS; 2) WATER, CARBON DIOXIDE}

\subsection*{TOSS-UP 17}
\textbf{MATH - Multiple Choice} The product of all the prime numbers between 10 and 20 lies in which of the following open intervals? \\
W) 10,000 to 20,000 \\
X) 20,000 to 40,000 \\
Y) 40,000 to 80,000 \\
Z) 80,000 to 160,000 \\
\textbf{ANSWER: Y) 40,000 TO 80,000}

\subsection*{VISUAL BONUS 17}
\textbf{MATH - Short Answer} Your next bonus is visual. Shown in this image are the first four iterations of a fractal. The original triangle has area 1. What is the fractal's final area? \\
\textbf{ANSWER: 8/5}

\subsection*{TOSS-UP 18}
\textbf{ENERGY - Multiple Choice} Researchers from the Shin Group explored using graphics processing units, or GPUs, to solve computationally intensive nonlinear programming problems. Which of the following tasks are GPUs most directly optimized to perform? \\
W) Frequent context switching \\
X) Deep recursion with branching \\
Y) Searching for large prime numbers \\
Z) Multiplying dense matrices \\
\textbf{ANSWER: Z) MULTIPLYING DENSE MATRICES}

\subsection*{VISUAL BONUS 18}
\textbf{ENERGY - Short Answer} Your next bonus is visual. Shown in the image is a linear classification problem studied by researchers at CSAIL. The dataset contains two distinct classes of white and black points, respectively. Answer the following two questions:
1) Identify all of the three hyperplanes $H_{1}$, $H_{2}$ and $H_{3}$ that separate the two classes;
2) Which one of the three hyperplanes would be chosen by a support vector machine? \\
\textbf{ANSWER: 1) $H_{2}$, $H_{3}$; 2) $H_{3}$}

\subsection*{TOSS-UP 19}
\textbf{PHYSICS - Short Answer} Lucas observes an atom with an vacancy in one of its inner electron shells suddenly ejecting an electron while another fills the vacancy. What is the name of this effect? \\
\textbf{ANSWER: AUGER EFFECT}

\subsection*{BONUS 19}
\textbf{PHYSICS - Short Answer} A system is in thermal equilibrium with a heat bath at temperature T, and has two states with equal degeneracies and energies $k_{B}T$ and $4k_{B}T$, respectively, where $k_{B}$ is the Boltzmann constant. To the nearest multiple of 5\%, what is the probability of a particle occupying the system's ground-state? \\
\textbf{ANSWER: 95\%}

\subsection*{TOSS-UP 20}
\textbf{CHEMISTRY - Short Answer} The Evans-Polanyi principle provides a relationship between a reaction's change in enthalpy and its activation energy. What hypothesis rationalizes this principle using the geometry of the reaction's transition state? \\
\textbf{ANSWER: HAMMOND POSTULATE}

\subsection*{BONUS 20}
\textbf{CHEMISTRY - Multiple Choice} The proton NMR spectrum of an organic compound contains a signal at 1.1 ppm with an integration of 6, and a second signal at 2.6 ppm with an integration of 1 that has been split into 7 peaks. Which of the following could be the chemical environment of the proton responsible for the second signal? \\
W) Proton adjacent to one aryl group \\
X) Proton adjacent to two aryl groups \\
Y) Proton adjacent to one methyl group \\
Z) Proton adjacent to two methyl groups \\
\textbf{ANSWER: Z) PROTON ADJACENT TO TWO METHYL GROUPS}

\subsection*{TOSS-UP 21}
\textbf{EARTH AND SPACE - Short Answer} Our Sun has absolute magnitude 4.8. A white dwarf is observed with absolute magnitude 14.8 and radius 0.01 solar radii. What is the ratio of the white dwarf's temperature to the Sun's temperature? \\
\textbf{ANSWER: 1}

\subsection*{BONUS 21}
\textbf{EARTH AND SPACE - Multiple Choice} A telescope's image of a point-like object often contain distinctive artifacts from that specific telescope. Which of the following telescopes is NOT properly matched with its artifacts of its images? \\
W) Hubble Space Telescope, four spokes \\
X) Spitzer Space Telescope, six equal spokes \\
Y) James Webb Space Telescope, six large spokes and two small spokes \\
Z) Chandra X-ray Observatory, three spokes \\
\textbf{ANSWER: Z) CHANDRA X-RAY OBSERVATORY, THREE SPOKES}

\subsection*{TOSS-UP 22}
\textbf{BIOLOGY - Short Answer} In order for Theca cells to convert testosterone to estrogen, they must produce what enzyme that converts androgens to estrogens? \\
\textbf{ANSWER: AROMATASE}

\subsection*{BONUS 22}
\textbf{BIOLOGY - Short Answer} GIP and GLP are both glucagon peptides that induce the release of insulin upon oral intake of food. They both belong to what more specific group of hormones, which are responsible for decreasing blood sugar? \\
\textbf{ANSWER: INCRETINS}

\subsection*{TOSS-UP 23}
\textbf{MATH - Multiple Choice} A geometric sequence has first term 4 and seventh term 9. For which of the following integers n is the nth term of this sequence rational? \\
W) 12 \\
X) 14 \\
Y) 16 \\
Z) 18 \\
\textbf{ANSWER: Y) 16}

\subsection*{BONUS 23}
\textbf{MATH - Short Answer} A rectangular field is 100 yards long and x yards wide. A horse travels at 28 miles per hour and a donkey at 20. Suppose the horse travels from one corner of the field to the other via edges, and the donkey goes diagonally. What is the minimum value of x such that the donkey takes at most as much time as the horse to finish this journey? \\
\textbf{ANSWER: 75}

\subsection*{TOSS-UP 24}
\textbf{ENERGY - Short Answer} The Van Voorhis Group is extending molecular dynamic simulations to long timescales where nuclei are no longer approximated to be stationary relative to electrons. What is this approximation known as? \\
\textbf{ANSWER: BORN-OPPENHEIMER}

\subsection*{BONUS 24}
\textbf{ENERGY - Short Answer} Physicists at the Laboratory for Nuclear Science supported by the DOE studied mirror nuclei of silicon isotopes. What nuclide is the mirror of silicon-32? \\
\textbf{ANSWER: ARGON-32}

\subsection*{TOSS-UP 25}
\textbf{BIOLOGY - Multiple Choice} Alex is using tandem mass spectrometry in his laboratory. Which of the following goals is he trying to accomplish? \\
W) DNA sequencing \\
X) Protein sequencing \\
Y) RNA sequencing \\
Z) Chromatin mapping \\
\textbf{ANSWER: X) PROTEIN SEQUENCING}

\subsection*{BONUS 25}
\textbf{BIOLOGY - Multiple Choice} Katya performs an experiment in which she isolates human mitochondria and places them into a solution of Tris-HCl buffered at pH 9.5. Given that the outer mitochondrial membrane is permable to protons, what will happen to the rate of ATP production, and the oxygen consumption of the cell, respectively? \\
W) Increase, Increase \\
X) Increase, Decrease \\
Y) Decrease, Increase \\
Z) Decrease, Decrease \\
\textbf{ANSWER: Y) DECREASE, INCREASE}

\end{document}
